\chapter{Why I Should Learn Math}
Millions of children hate math, and often for good reasons: the content feels
dry, the teaching feels lifeless, and many successful people seem to live well
without mathematics.

So let us ask the question honestly: why should I learn math?  Why sit in class
learning \(1+1\) when watching a movie feels far more enjoyable?

\section{Bad Answers}

I asked AI, and it gave me the usual textbook answers.  Below I separate each
\textbf{Answer} from \textbf{My critique}.

\subsection*{Answer 1}
\begin{quote}
``Learning math is essential because it enhances problem-solving skills,
supports daily life decisions, strengthens cognitive abilities, and fosters
creativity.  Practical applications in daily life include cooking, budgeting,
discounts, time management, travel planning, and home projects.''
\end{quote}

\subsection*{My Critique 1}
Come on.  A great chef does not stand in the kitchen doing explicit
calculations all day.  He or she knows what to do and when to do it, often by
experience and judgment.

\subsection*{Answer 2}
\begin{quote}
``Studying math strengthens your brain by improving memory, logical reasoning,
and analytical thinking.  Math problem-solving teaches you to identify relevant
information and apply systematic strategies that transfer to life and work.''
\end{quote}

\subsection*{My Critique 2}
I do not buy this when it is presented as a slogan.  Too often it sounds like
an unmotivated, unhappy teacher trying to convince students of something he
does not fully believe himself.

\subsection*{Answer 3}
\begin{quote}
``Math is foundational for many careers, especially STEM, finance, and
technology.  It also builds precision, objectivity, critical thinking, and
persistence, which help in almost any profession.''
\end{quote}

\subsection*{My Critique 3}
There are many rich people -- soccer players, actors, singers -- who are not
known for being strong at math.  So ``learn math to become successful'' sounds
unconvincing if success is measured only by money or fame.

\section{My Friend's Answer}
A friend of mine runs a startup and has four children.  One day, his
first-grade daughter Luisa asked why she needed to learn math.

At first he gave the usual abstract answer: ``As a business owner, I need to
understand income and expenses.''  Then he made it concrete: ``How about we
bake some cakes, sell them, and calculate profit and loss together?''

They bought ingredients, baked the cakes, and brought them first to the school
bus and then to school.

It became a small sensation.  Even the school drivers loved the cakes, and they
became a topic of conversation for the whole day.

I do not know whether Luisa fully understood the profit-and-loss calculation.
But she clearly had fun, and the math finally had a real context.

\section{My Answer}

\subsection*{Case 1: Ed Thorp and Blackjack}
Ed Thorp did not come from wealth.  He used mathematics to attack a concrete
question: what are the true odds in blackjack if cards are not replaced
immediately?  By computing conditional probabilities and expected value under
different deck compositions, he showed that the game can shift from house edge
to player edge in specific states.

To scale the strategy
into real money, he needed outside capital; one famous story is that he even
used newspaper ads to find backers, and one sponsor reportedly had underworld
connections. That was not just theory.  Casinos reacted.  He and his teammates had to use
disguises and operational discipline to avoid detection.  

The point is not glamour.  The point is precision: math turned ``I feel lucky''
into a quantified edge.

\subsection*{Case 2: LTCM and Model Risk}
Long-Term Capital Management (LTCM) had elite people and sophisticated models.
They were not ignorant of mathematics; they relied on it heavily.  But they
treated their model assumptions as if they were stable laws of nature.

Their positions were built on spreads that looked historically safe and mean
reverting.  When market conditions changed sharply, correlations moved, liquidity
dried up, and losses exploded faster than their risk framework expected.  The
fund nearly collapsed the financial system and required a coordinated rescue.

This is the second lesson: bad math, or good math used with false assumptions,
can destroy you.  Learning math is not only about finding opportunities; it is
also about understanding model limits, tail risk, and what can go wrong when
the world refuses to follow your equations.

\subsection*{Case 3: Voltaire and the French Lottery}
Historically, Voltaire (Fran\c{c}ois-Marie Arouet, 1694--1778) was one of the
central figures of the French Enlightenment: a prolific writer, satirist, and
public intellectual known for works such as \emph{Candide}, and for his
criticism of intolerance and abuse of power.

Voltaire is remembered as a writer and philosopher, but he also became wealthy
through a mathematically informed lottery strategy in France.  He and his
partners noticed a pricing flaw in the state lottery: under certain structures
of ticket pricing and payouts, buying enough tickets could create a positive
expected value.

This was not blind gambling.  It was arithmetic: identify a mispriced system,
size the participation, and repeat while the edge exists.  Voltaire's fortune
was built less by luck than by exploiting a structural inconsistency that most
people did not quantify.

\subsection*{Case 4: Casanova and the Same Lottery}
Giacomo Casanova (1725--1798), according to memoirs and historical accounts,
was not only a famous libertine adventurer but also a diplomat, writer,
occasional spy, and social operator who moved through courts and salons across
Europe.  He became widely known for his autobiography \emph{Histoire de ma vie}
(\emph{The Story of My Life}).

Casanova, in contrast, made money by helping promote and sell the same French
lottery.  He understood the social psychology around it and turned that into
wealth.

But later he was ruined by gambling losses.  The same man who benefited from a
system with engineered odds lost control when facing games where the odds were
against him.  The practical lesson is sharp: understanding one profitable setup
does not protect you if you misread probability in another.

\section{Conclusion}
You do not need to learn math to live a great life.  Running a fund like LTCM
and blowing up, or escaping prison in Venice and later going broke at gambling
tables, is also a kind of great life story.

So the choice is yours.  Do you want to be simple and rich like Voltaire and
Thorp, or spectacular and broke like LTCM or Casanova?